\begin{lexlean}{DevRel}
\useglossary{lexlean.std.bool@1.2.0}
\useglossary{lexlean.std.nat@1.2.0}
\useglossary{prism.arch@1.0.0}
\importmodule{Catalog}
\importmodule{Foundation.DevRel.V1.Model}
\importmodule{Foundation.View.Text.V1.Model}
\importmodule{Foundation.View.Text.V3.Model}
\title{DevRel record lookup}
\begin{semanticmodule}
definition{
  name: dispatch,
  parameters: [value : Bytes],
  result: Bytes,
  body: (Foundation.DevRel.V1.Model::lookup value (Catalog::records) (Catalog::resources) (Catalog::covers)),
  axioms: [Quot.sound, propext]
}

definition{
  name: application,
  parameters: [],
  result: Foundation.View.Text.V3.Model::TextApplication,
  body: (record Foundation.View.Text.V3.Model::TextApplication {profile := "prismpm/text-application/3", name := "Prism-DevRel", cargoName := "prism-devrel", cargoVersion := "0.1.0", cargoDescription := "Answers which learning resources teach a facet record of the subject Prism model.", cargoRepository := "https://github.com/afflom/prism-devrel", cargoHomepage := "https://afflom.github.io/prism-devrel/", libraryRoots := (cons "PrismDevRel.DevRel.dispatch" (nil String)), acceptanceVectors := (Catalog::vectors), entryRoot := "PrismDevRel.DevRel.dispatch", coreContract := "hologram:guest/core-wasm@1", requestMaximum := (16 : UInt32), responseMaximum := (1024 : UInt32), guestAllocationMaximum := (17 : UInt32), coreMemoryPages := (ctor Option.some [UInt32] (32 : UInt32)), coreStackBytes := (ctor Option.some [UInt32] (65536 : UInt32)), capabilitiesEmpty := true, fatArchive := true, primaryLayer := (0 : UInt8), viewLayer := (1 : UInt8), view := (record Foundation.View.Text.V1.Model::TextView {title := "Prism-DevRel", heading := "Prism-DevRel", inputLabel := "Record index", submitLabel := "Look up", outputLabel := "Record and resources", inputError := "Enter a record index of at most six digits.", responseError := "The answer could not be displayed."}), cli := (record Foundation.View.Text.V3.Model::TextCli {command := "lookup", description := "Prints the identifier of the subject record at an index and the slugs of the learning resources that cover it.\n", usageCode := (64 : UInt8), ioCode := (74 : UInt8), limitCode := (ctor Option.none [UInt8]), outcomes := (cons (record Foundation.View.Text.V3.Model::TextCliOutcome {name := "covered", prefix := "record ", exitCode := (0 : UInt8)}) (cons (record Foundation.View.Text.V3.Model::TextCliOutcome {name := "unknown", prefix := "unknown record", exitCode := (1 : UInt8)}) (cons (record Foundation.View.Text.V3.Model::TextCliOutcome {name := "error", prefix := "invalid request", exitCode := (2 : UInt8)}) (nil Foundation.View.Text.V3.Model::TextCliOutcome)))), sections := (cons (record Foundation.View.Text.V3.Model::TextCliSection {title := "REQUEST", body := "A record index in decimal, one to six digits.\n"}) (nil Foundation.View.Text.V3.Model::TextCliSection)), examples := (cons (record Foundation.View.Text.V3.Model::TextCliExample {name := "first", text := "0"}) (nil Foundation.View.Text.V3.Model::TextCliExample))}), workbench := (ctor Option.none [Foundation.View.Text.V3.Model::WorkbenchView])})
}
\end{semanticmodule}
\end{lexlean}
